\chapter{How to Read This Book}
\label{ch:how_to_read}

This book defines every finance term on first use, derives every result
that teaches something, and stops short of derivations that do not. It is
written for a strong mathematics and computer science student with
\emph{zero} prior exposure to finance. Existing resources either assume
the reader already knows what a bid--ask spread is, or dilute the
mathematics to the point where nothing can be derived; this book sits in
the gap.

Two concrete goals share a deadline in mid-April 2026. \textbf{IMC
Prosperity~4} is a two-week algorithmic trading competition starting
2026-04-14 in which teams write Python algorithms to trade simulated
products against each other and scripted bots. Past editions rewarded
market making, basket statistical arbitrage, options pricing, locational
arbitrage, and exploitation of bot behaviour --- all taught here from
first principles. The \textbf{Goldman Sachs Quant Strats off-cycle
placement} begins around 2026-04-20; the reader will sit on a trading
desk as a strategist (``Strats''). Since the desk assignment is not yet
known, Part~VII covers every major asset class --- rates, equities, FX,
credit, commodities --- at comparable depth, together with production
risk, numerical methods, and firm-specific infrastructure (SecDB, Slang).

Parts~I--V form the common core; Part~VI is Prosperity-specific and
Part~VII Goldman-specific. A reader pursuing only one goal can skip the
other Part entirely. The reading paths in
Section~\ref{sec:reading_paths} make this concrete.

\section{Who this book is for}
\label{sec:who_for}

The reader is assumed to have the following background.

\paragraph{Mathematics.} Undergraduate probability (random variables,
expectation, conditional expectation, the central limit theorem, basic
hypothesis testing), linear algebra (vector spaces, eigenvectors, SVD at
the mechanical level), multivariable calculus, and some exposure to
stochastic processes at the level of Brownian motion, It\^o's lemma, and
simple stochastic differential equations. Measure theory is \emph{not}
assumed; where a measure-theoretic statement would be cleaner, the book
states it informally and cites the rigorous version. Basic optimisation
(Lagrange multipliers, convex objectives) is used without reintroduction.

\paragraph{Computer science.} Python fluency, including \texttt{numpy},
\texttt{pandas}, and standard data structures; time complexity; basic
algorithms (sorting, hashing, dynamic programming); and the ability to
read a moderately-sized Python codebase unaided. Code in this book is
sparing --- at most one reference skeleton per strategy, typically
15--30 lines, pedagogical rather than production-ready. Full runnable
code lives outside the PDF.

\paragraph{Finance.} Zero. The reader is not assumed to know what a stock
is, what ``going long'' means, what an exchange does, what an option is,
or what the terms ``bid'', ``ask'', ``mid'', ``spread'', ``tick'',
``lot'', ``maker'', ``taker'', ``slippage'', ``delta'', ``vega'', or
``hedge'' mean. Each is defined on first use; no finance jargon appears
without a prior in-text definition. Appendix~B collects the master
notation table and a glossary of trading terms as a reference backstop.

Readers whose probability, linear algebra, or stochastic-calculus
background is rusty should consult Appendix~A before Part~II
(microstructure) and again before Part~IV (options). Appendix~A is a
refresher, not a course: it restates the results used and points to
standard sources for proofs. If the refresher is insufficient, spend a
week with a probability textbook before continuing.

\section{Three reading paths}
\label{sec:reading_paths}

The book has 39 numbered chapters across seven Parts. The three paths
below are each \emph{complete and self-contained} for their target use
case: a chapter not on a path is safe to skip for that goal; a chapter
on a path is load-bearing.

\subsection*{Path A: Prosperity speed-run}

\[
\begin{gathered}
\text{Ch}\ 1 \to 2 \to 3 \to 8 \to 9 \to 11 \to 14 \to 15 \to 16 \\
\to 23 \to 24 \to 25 \to 26.
\end{gathered}
\]

For a reader whose sole immediate goal is Prosperity~4. It covers
exchanges, participants, and price (Ch~1--3), then the minimum strategy
vocabulary the competition historically rewards: market making (Ch~8),
mean reversion (Ch~9), and basket statistical arbitrage (Ch~11). Enough
options theory to price and delta-hedge the voucher products follows
(Ch~14--16), then the full Prosperity playbook (Ch~23--26) --- tooling,
algorithmic round strategies, manual rounds, and lessons from top-team
writeups. Omitted: Parts~V and~VII, most of the arbitrage zoo beyond
baskets, and the full vol surface (Ch~17). Roughly 130 pages; intended
for a reader starting a week before Prosperity~4 begins.

\subsection*{Path B: Goldman Sachs Strats prep}

\[
\begin{gathered}
\text{Ch}\ 1 \to 2 \to 3 \to 14 \to 15 \to 16 \to 17 \to 27 \to 28 \to 29 \\
\to 30 \to 31 \to 32 \to 33 \to 34 \to 35 \to 36 \to 37 \to 38.
\end{gathered}
\]

For a reader showing up to a sell-side quant strats desk ready to
contribute. It begins with the minimum market-structure vocabulary
(Ch~1--3), then the full options stack (Ch~14--17) since every Part~VII
chapter cross-references it, then every Part~VII chapter in order: the
role itself (Ch~27), fixed income foundations --- the largest knowledge
gap for a CS reader (Ch~28), rates derivatives and SABR (Ch~29), equity
derivatives and variance swaps (Ch~30), FX derivatives (Ch~31), credit
derivatives (Ch~32), commodities (Ch~33), numerical methods (Ch~34),
production risk and P\&L Explain (Ch~35), SecDB and Slang (Ch~36),
regulatory context (Ch~37), and interview-grade problems (Ch~38).
Omitted: Parts~III, V, and VI, which teach buy-side strategy
construction --- not the Strats job. Roughly 500 pages; intended for a
reader with three to five weeks until the placement starts.

\subsection*{Path C: Full cover-to-cover}

All 39 chapters in order, plus Appendices~A--C. Intended for a reader
building a permanent foundation rather than preparing for a deadline.
Estimated effort: around 300 hours of focused reading, working every
derivation and worked example by hand. Hard-flagged chapters (see
Section~\ref{sec:hard_chapters}) should be budgeted at roughly double
the time of an average chapter.

\section{Chapter dependency map}
\label{sec:dep_map}

Any Part can be read after Parts~I and~II, but not before them.
Figure~\ref{fig:dep_map} shows the coarse dependency structure; the
prose below explains it.

\begin{figure}[htbp]
\centering
\begin{tikzpicture}[
  node distance=8mm and 14mm,
  part/.style={draw, rounded corners, align=center, minimum height=9mm,
               minimum width=22mm, font=\small},
  arrow/.style={-{Latex[length=2mm]}, thick}
]
  \node[part, fill=gray!10] (P1) {Part I\\Foundations};
  \node[part, fill=gray!10, right=of P1] (P2) {Part II\\Microstructure};
  \node[part, fill=gray!10, right=of P2] (P3) {Part III\\Strategies};
  \node[part, fill=gray!10, below=of P2] (P4) {Part IV\\Options};
  \node[part, fill=gray!10, right=of P4] (P5) {Part V\\Execution / Risk};
  \node[part, fill=gray!10, below=of P4] (P6) {Part VI\\Prosperity};
  \node[part, fill=gray!10, right=of P6] (P7) {Part VII\\GS Strats};

  \draw[arrow] (P1) -- (P2);
  \draw[arrow] (P2) -- (P3);
  \draw[arrow] (P1) -- (P4);
  \draw[arrow] (P3) -- (P5);
  \draw[arrow] (P4) -- (P5);
  \draw[arrow] (P3) -- (P6);
  \draw[arrow] (P4) -- (P6);
  \draw[arrow] (P4) -- (P7);
  \draw[arrow] (P5) -- (P7);
\end{tikzpicture}
\caption{Coarse dependency graph between Parts. An arrow $X \to Y$ means
$Y$ assumes the content of $X$. Parts~I and~II are prerequisites for
everything; Part~IV (options) is a prerequisite for both Part~VI and
Part~VII; Part~V cross-references Parts~III and~IV.}
\label{fig:dep_map}
\end{figure}

At chapter granularity: Part~I (Ch~1--3) introduces exchanges, the
limit order book, market participants, and the many definitions of
``the price''; every later chapter assumes this vocabulary. Part~II
(Ch~4--7) builds the theory of \emph{why} spreads exist, underpinning
every Part~III strategy --- Ch~8 (market making) in particular makes no
sense without the adverse-selection argument of Ch~5. Inside Part~IV,
Ch~14 (payoffs, no-arbitrage bounds, put--call parity) precedes Ch~15
(Black--Scholes), which precedes Ch~16 (Greeks and delta hedging) and
Ch~17 (the volatility surface and gamma scalping). Part~VII assumes the
Part~IV options stack in full, then layers production, sell-side, and
asset-class-specific content on top: Ch~29 (rates derivatives)
cross-references Ch~15--17, Ch~30 extends Ch~17 with variance swaps,
and Ch~34 (numerical methods) is where the pricing problems of
Ch~29--32 get solved in practice. Part~V (execution and risk) sits
between the strategy and options material and the two target-specific
Parts: it covers how trades are actually placed without moving the
market against oneself, and how positions are sized and monitored so a
losing day does not become a career-ending one. Part~VI (Prosperity) is
an applied layer over Parts~I--IV rather than a prerequisite for
anything --- Part~VI is the most self-contained Part in the book.

\section{Notation}
\label{sec:notation}

Every symbol used across multiple chapters is defined once and
catalogued in Appendix~B, alongside a glossary of trading terms.
Appendix~B is a reference --- consult it whenever a symbol introduced
many chapters earlier needs refreshing. No chapter redefines a symbol
that appears in the master table, so a symbol's meaning is fixed the
first time it is used.

A few symbols are introduced locally within a single chapter (for
example, the half-life in the Ornstein--Uhlenbeck chapter, or the
hazard rate in the credit chapter); these do not recur globally and
are omitted from the master table.

\section{Callout boxes}
\label{sec:callouts}

Seven coloured callout boxes recur throughout the book, introduced
here so that a reader meeting one in Chapter~4 already knows what it
signals. This chapter uses no boxes --- they are reserved for technical
content.

\begin{itemize}
\item \textbf{Background} (gray). Prerequisite reminders at the start
of most chapters, listing the earlier-chapter results assumed. Skim if
confident, or follow the back-references if not.

\item \textbf{Key fact} (blue). A formal result worth memorising ---
the bid--ask spread formula, Kyle's $\lamk$, Avellaneda--Stoikov's
reservation-price formula. Every major result appears in a key-fact
box near where it is derived.

\item \textbf{Intuition} (green). A plain-English restatement of what a
formal result \emph{means}, placed immediately after the derivation.
When the reader has followed the algebra but is unsure what was
shown, the intuition box is the answer. \emph{Slow down and read these
carefully.}

\item \textbf{Prosperity example} (orange). A worked example using a
specific Prosperity product from a prior edition. Appears inline in
Parts~I--V wherever theory maps cleanly onto a Prosperity product. The
full playbook is Part~VI.

\item \textbf{On a Goldman Sachs desk} (teal). How the technique just
taught is used on a sell-side trading desk: what the Strat writes,
which system owns the calculation, and which desk cares most. Appears
mostly in Parts~IV and~V; Part~VII treats the same material in depth.

\item \textbf{Common mistake / Pitfall} (red). A wrong intuition that
sounds right, or a failure mode the reader will otherwise walk into.
\emph{Slow down and read these carefully.} These boxes are the highest
value-per-word in the book --- they encode mistakes that cost
professionals money.

\item \textbf{Historical note} (purple). An aside on the origin of a
result, a biographical note, or a short story about how a model came
to dominate practice. Skip on a first read without losing anything
technical.
\end{itemize}

\section{How to read a hard chapter}
\label{sec:hard_chapters}

Five chapters are flagged as \emph{hard}: Ch~5 (Glosten--Milgrom and
Kyle), Ch~8 (Avellaneda--Stoikov market making), Ch~15 (Black--Scholes
from first principles), Ch~22 (the ``which strategy when'' decision
framework), and Ch~29 (SABR and the rates vol surface). Hard means the
chapter contains a derivation or synthesis which, on first reading,
feels dense regardless of the reader's mathematical background. These
deserve extra time and, usually, a second pass.

Every hard chapter follows the same four-stage pattern, mirroring the
worked-example style of the NLP lecture-notes reference this book
takes its form from.

\begin{enumerate}
\item \textbf{Toy example first.} The chapter opens with a concrete
numerical example worked by hand with small integers, before any
symbol is defined abstractly --- so the reader experiences the
phenomenon before formalising it. Ch~5, for instance, opens with a
Glosten--Milgrom game on two states $\{V=100, V=110\}$ with known
probabilities; the reader computes the bid and ask by hand before
seeing the general formula.

\item \textbf{Formalise.} The same problem is restated in general
notation. No new results are proved here; the goal is to translate
the toy example into the language the derivation will use.

\item \textbf{Derive.} The main result is derived line by line.
Routine steps are flagged and skipped; illuminating ones are written
out in full with underbraces labelling each term. Derivations follow
rigour level~B: rigorous where the algebra teaches something,
stated-with-citation where it does not.

\item \textbf{Intuition box.} The chapter closes with an intuition box
restating the result in one or two sentences of plain English, often
with a picture. If the reader got lost in the derivation, the intuition
box is the minimum they need to carry forward.
\end{enumerate}

Recommended reading protocol: work stage~1 carefully with a pen,
checking every number in the toy example is reproducible from the
rules. Read stages~2 and~3 once at normal speed, then re-read while
writing down each line of algebra on paper. Finally, read the intuition
box and ask whether it would have been obvious from the toy example
alone --- if yes, the chapter has done its job. If no, re-read
stages~3 and~4 together and look for what the intuition box points at
that the derivation glossed over.

Budget roughly double the time for a hard chapter versus an average
chapter in the same Part. The five hard chapters contain most of the
load-bearing theory; time spent on them compounds because every later
chapter assumes the result without rederiving it.
